\documentclass[11pt]{amsart}
\usepackage[T1]{fontenc}
\usepackage{amsmath,amssymb,amsthm}
\usepackage{hyperref}
\usepackage{enumitem}

\newtheorem{theorem}{Theorem}
\newtheorem{lemma}[theorem]{Lemma}
\newtheorem{corollary}[theorem]{Corollary}
\theoremstyle{definition}
\newtheorem{definition}[theorem]{Definition}
\theoremstyle{remark}
\newtheorem{remark}[theorem]{Remark}

\newcommand{\X}{X}
\newcommand{\T}{\mathcal T}
\newcommand{\B}{B_X}
\newcommand{\intt}{\operatorname{int}}
\newcommand{\diam}{\operatorname{diam}}

\title{A Boundary-Singularity Obstruction for Rotund Ball Tilings}
\author{agent\_lane\_19}
\date{June 19, 2026}

\begin{document}
\maketitle

\section*{Source Problem}

The source paper is Carlo Alberto De Bernardi, Tommaso Russo, and Jacopo
Somaglia, \emph{Lattice tilings of Hilbert spaces}, arXiv:2505.04267. In
Problem 7.5 the authors ask:

\begin{quote}
Does there exist a rotund and/or smooth normed space \(X\) that admits a tiling
by balls?
\end{quote}

The source paper recalls that separable spaces do not admit tilings by rotund
or smooth bounded convex bodies, while in the nonseparable case the known
negative results cover LUR or Frechet smooth spaces. The result below is a
partial obstruction for arbitrary rotund ball tilings. The final corollary
records a diametral strengthening obtained in a later push on the same
problem.

\section*{Definitions}

A family \(\T\) of bodies tiles a normed space \(X\) if it covers \(X\) and the
members have pairwise disjoint interiors. A point \(x\in X\) is regular for
\(\T\) if some neighbourhood of \(x\) meets only finitely many members of
\(\T\); otherwise it is singular.

A boundary point \(x\) of a convex body \(B\) is a quasi-polyhedral point
(QP-point) if there is a neighbourhood \(U\) of \(x\) such that
\([x,y]\subset \partial B\) whenever \(y\in U\cap\partial B\). A norm is rotund
if its unit sphere contains no non-trivial segment.

A point \(x\in S_X\) is locally non-D2 (LND2) if there is \(\delta>0\) such
that
\[
  \diam\{y\in S_X:\|(x+y)/2\|\geq 1-\delta\}<2.
\]
This condition is weaker than local uniform rotundity at \(x\), but it still
rules out diameter-two behaviour in arbitrarily thin midpoint neighbourhoods
of \(x\).

\section*{Idea of the Proof}

Regularity of a boundary point is a local finiteness condition. Near such a
point only finitely many convex tiles are visible, so the tile containing the
point is locally cut out by finitely many supporting half-spaces. This is
exactly the quasi-polyhedral mechanism used in the source paper. But a
quasi-polyhedral point on the boundary of a body in dimension at least two
forces a small boundary segment. Therefore a rotund unit ball can never have a
regular boundary point in a ball tiling.

The strengthened obstruction uses the opposite case. If a boundary point is
singular, then arbitrarily close to it one sees two further unit balls. A
standard three-balls lemma converts these three nearby pairwise non-overlapping
balls into pairs of unit vectors whose midpoints stay arbitrarily close to the
sphere while their mutual distance is arbitrarily close to \(2\). Thus
singularity forces failure of LND2.

\begin{lemma}[regular boundary points are QP]\label{lem:regular-qp}
Let \(\T\) be a convex tiling of a normed space \(X\). Let \(B\in\T\), and let
\(x\in\partial B\) be regular for \(\T\). Then \(x\) is a QP-point of \(B\).
\end{lemma}

\begin{proof}
Since \(x\) is regular, there is a neighbourhood \(U_0\) of \(x\) that meets
only finitely many tiles. Let these tiles be
\[
  B,B_1,\ldots,B_n.
\]
Shrinking \(U_0\), we may assume that
\[
  U_0\subset \intt(B\cup B_1\cup\cdots\cup B_n).
\]
For each \(k\), the convex bodies \(B\) and \(B_k\) have disjoint interiors and
meet at \(x\), so the Hahn--Banach separation theorem gives a functional
\(f_k\in X^*\) such that
\[
  \sup f_k(B)=f_k(x)=\inf f_k(B_k).
\]
After possibly shrinking to a convex neighbourhood \(U\subset U_0\), the local
description
\[
  B\cap U
  =
  U\cap\bigcap_{k=1}^n\{z:f_k(z)\leq f_k(x)\}
\]
holds. Indeed, if \(z\in U\setminus B\), then \(z\) lies in one of the adjacent
tiles \(B_k\), and the corresponding separating functional excludes \(z\) from
the right-hand side.

Thus \(x\) is locally a finite-half-space boundary point, in particular a
QP-point. This is the same local observation recorded as Fact 3.8 in the
source paper.
\end{proof}

\begin{lemma}[QP points obstruct rotundity]\label{lem:qp-rotund}
Let \(B\) be the closed unit ball of a normed space \(X\) with
\(\dim X\geq 2\). If \(\partial B\) has a QP-point, then \(B\) is not rotund.
\end{lemma}

\begin{proof}
Let \(x\in\partial B\) be a QP-point and choose the neighbourhood \(U\) from
the definition. Since \(\dim X\geq2\) and \(B\) is a convex body, \(\partial B\)
has no isolated points. Hence there is \(y\in U\cap\partial B\) with
\(y\neq x\). By the QP property, the non-trivial segment \([x,y]\) is contained
in \(\partial B\). Therefore the unit sphere contains a non-trivial segment,
so \(B\) is not rotund.
\end{proof}

\begin{theorem}[boundary-singularity obstruction]\label{thm:main}
Let \(X\) be a normed space with \(\dim X\geq2\). Suppose that
\[
  \T=\{a_i+\B:i\in I\}
\]
is a tiling of \(X\) by translates of the closed unit ball. If the norm of
\(X\) is rotund, then every boundary point of every tile is singular for
\(\T\).
\end{theorem}

\begin{proof}
It suffices to consider the tile \(\B\), since translations preserve both
regularity and rotundity. Suppose that \(x\in\partial\B\) is regular for
\(\T\). By Lemma~\ref{lem:regular-qp}, \(x\) is a QP-point of \(\B\). By
Lemma~\ref{lem:qp-rotund}, \(\B\) is not rotund, contradicting the assumption.
Thus no boundary point can be regular.
\end{proof}

\begin{corollary}
No normed space of dimension at least two with a rotund norm admits a locally
finite tiling by translates of its closed unit ball.
\end{corollary}

\begin{proof}
In a locally finite tiling every point is regular. This contradicts
Theorem~\ref{thm:main}, because the unit ball has nonempty boundary.
\end{proof}

\begin{lemma}[three nearby balls force local diameter two]\label{lem:three-balls}
Let \(X\) be a normed space, let \(B_0=B_X\), and let \(B_1,B_2\) be closed
balls of radius one whose interiors are pairwise disjoint from each other and
from \(\intt B_0\). Suppose that \(y_i\in\partial B_i\) for \(i=0,1,2\), and
that \(\diam\{y_0,y_1,y_2\}\leq \eta\). Then
\[
  \diam\{y\in S_X:\|y_0+y\|\geq 2-\eta\}\geq 2-2\eta .
\]
\end{lemma}

\begin{proof}
This is the standard three-balls estimate of De Bernardi--Somaglia--Vesely
\cite[Lemma 3.7]{DSV} and De Bernardi--Vesely \cite[Lemma 4.5]{DV}. It is
also quoted in the packing-constant literature in exactly this form. We use it
only as a local geometric lemma.
\end{proof}

\begin{lemma}[singular boundary points fail LND2]\label{lem:singular-lnd2}
Let \(\T=\{a_i+B_X:i\in I\}\) be a tiling by translates of the closed unit ball.
If \(x\in S_X\) is singular for \(\T\), then for every \(\delta>0\),
\[
  \diam\{y\in S_X:\|(x+y)/2\|\geq 1-\delta\}=2.
\]
In particular, \(x\) is not an LND2 point.
\end{lemma}

\begin{proof}
Fix \(\eta>0\). Since \(x\) is singular, every neighbourhood of \(x\) meets
infinitely many tiles. Choose two distinct non-central tiles \(T_1,T_2\) that
meet \(B(x,\eta)\). For \(i=1,2\), choose \(z_i\in T_i\cap B(x,\eta)\).

The point \(x\) is not an interior point of \(T_i\): otherwise a small
neighbourhood of \(x\) inside \(T_i\) would meet \(\intt B_X\), contradicting
the disjointness of tile interiors. Hence the segment from \(z_i\) to \(x\)
contains a point \(y_i\in\partial T_i\) with \(\|y_i-x\|\leq \eta\). Taking
\(y_0=x\in\partial B_X\), the three unit balls \(B_X,T_1,T_2\) satisfy the
hypotheses of Lemma~\ref{lem:three-balls} with
\(\diam\{y_0,y_1,y_2\}\leq 2\eta\). Therefore
\[
  \diam\{y\in S_X:\|x+y\|\geq 2-2\eta\}\geq 2-4\eta .
\]
Equivalently,
\[
  \diam\{y\in S_X:\|(x+y)/2\|\geq 1-\eta\}\geq 2-4\eta .
\]
For a fixed \(\delta>0\), let \(\eta<\delta\) tend to \(0\). The displayed
sets increase when the threshold is weakened, so the set corresponding to
\(\delta\) has diameter \(2\). Thus \(x\) is not LND2.
\end{proof}

\begin{corollary}[diametral strengthening]\label{cor:no-lnd2}
Let \(X\) be a normed space with \(\dim X\geq2\). If translates of \(B_X\) tile
\(X\) and the norm of \(X\) is rotund, then no point of \(S_X\) is LND2.
\end{corollary}

\begin{proof}
By Theorem~\ref{thm:main}, every point of \(S_X=\partial B_X\) is singular for
the tiling. Lemma~\ref{lem:singular-lnd2} then applies at every point of the
unit sphere.
\end{proof}

\begin{corollary}[smooth no-flat variant]
Let \(X\) admit a tiling by translates of its closed unit ball. If a boundary
point \(x\) of a tile is regular and the norm is Gateaux smooth at \(x\), then
\(x\) is a flat point of the unit ball. Consequently, a Gateaux smooth ball
tiling with no flat boundary points must have every boundary point singular.
\end{corollary}

\begin{proof}
The regular point is QP by Lemma~\ref{lem:regular-qp}. The standard local
convexity observation recalled in the source paper says that every Gateaux
smooth QP-point is flat: the unique supporting hyperplane at a smooth
quasi-polyhedral point must contain a relative neighbourhood of the boundary.
The final statement is immediate.
\end{proof}

\section*{Limitations and Novelty Check}

This is a partial obstruction, not a solution to Problem 7.5. It leaves open
whether a rotund or smooth normed space with a necessarily everywhere-singular
ball tiling exists. The strengthened conclusion shows that a positive rotund
example would have to be extreme in a second way: every unit vector must fail
LND2, so the unit sphere must be rotund but locally diameter-two at every
point.

I checked the run indexes for arXiv:2505.04267 and the phrases
\emph{lattice tilings}, \emph{rotund}, \emph{smooth}, and \emph{ball tiling};
no duplicate packet appeared. A bounded web search found the source arXiv page
and no later exact resolution of Problem 7.5. The argument intentionally uses
the source paper's local QP facts, so human review should compare the statement
against the older Klee--Maluta--Zanco and Klee--Tricot literature.

\begin{thebibliography}{9}

\bibitem{DRS}
C.~A. De Bernardi, T. Russo, and J. Somaglia,
\emph{Lattice tilings of Hilbert spaces}, arXiv:2505.04267.

\bibitem{DSV}
C.~A. De Bernardi, J. Somaglia, and L. Vesely,
\emph{Star-finite coverings of Banach spaces},
J. Math. Anal. Appl. 491 (2020), 124384.

\bibitem{DV}
C.~A. De Bernardi and L. Vesely,
\emph{Tilings of normed spaces},
Canad. J. Math. 69 (2017), 321--337.

\bibitem{KMZ}
V. Klee, E. Maluta, and C. Zanco,
\emph{Tiling with smooth and rotund tiles},
Fund. Math. 126 (1986), 269--290.

\bibitem{KT}
V. Klee and C. Tricot,
\emph{Locally countable plump tilings are flat},
Math. Ann. 277 (1987), 315--325.

\end{thebibliography}

\end{document}
